\rhead{MN-CIV: Urban Informatics \& Smart Infrastructure}
\phantomsection
\section*{Urban Informatics \& Smart Infrastructure (MN-CIV)}
\label{minor:CIV_L1}

\noindent \small{\hyperref[main:scheme]{$\leftarrow$ Back to Scheme}} \vspace{0.5cm}

\noindent \textbf{Credits:} 2 (2-0-0)

\subsection*{Course Content}
\noindent\textbf{Unit 1} \\
\textit{The City as a Computational System:} Urban informatics defined: the application of data science, sensing, and computation to understand and manage cities as complex adaptive systems; The smart city stack: physical infrastructure $\rightarrow$ sensor and actuator layer $\rightarrow$ communication networks $\rightarrow$ data platforms $\rightarrow$ analytics and applications as a layered architecture; Urban data taxonomy: administrative records (census, permits), transactional data (transit ticketing, utility billing), sensor streams (traffic loops, air quality monitors), and participatory data (citizen reports, social media) as heterogeneous sources with different latency and quality profiles; City-scale data platforms: urban operating system (Urban OS) architectures and open data portals as the infrastructure for data governance and third-party application development; The digital divide and data equity: uneven sensor deployment, demographic underrepresentation in datasets, and the risk of algorithmic systems reinforcing spatial inequality; Key urban indicators: population density, land use mix, walkability index, and heat island intensity as computed metrics derived from raw urban data.

\vspace{0.3cm}

\noindent\textbf{Unit 2} \\
\textit{Urban Sensing Networks and Data Acquisition:} Fixed sensing infrastructure: inductive loop detectors, pneumatic tube counters, CCTV-based vehicle detection, and environmental sensor networks as the permanent data collection layer of a city; Mobile sensing: probe vehicles (taxis, buses with GPS), crowdsourced pothole detection via accelerometers, and participatory sensing apps as opportunistic data collection paradigms; Communication infrastructure for smart cities: LoRaWAN for low-power wide-area sensor networks, NB-IoT for deep-indoor metering, 5G for latency-sensitive applications (V2X, remote surgery), and DSRC/C-V2X for vehicle-to-infrastructure communication; Smart metering for utilities: AMI (Advanced Metering Infrastructure) for electricity, water, and gas; interval data as a high-frequency time series enabling disaggregation of end-use consumption; Data fusion from heterogeneous sources: temporal alignment, spatial co-registration, and handling missing data in multi-sensor urban streams; Edge-cloud continuum: processing latency-sensitive urban data locally (traffic signal control) vs.\ aggregating historical data centrally (urban planning analytics).

\vspace{0.3cm}

\noindent\textbf{Unit 3} \\
\textit{Urban Mobility and Transportation Systems:} Transportation networks as weighted directed graphs: nodes as intersections, edges as road segments with travel time, capacity, and incident attributes; Traffic flow theory: the fundamental diagram relating flow, density, and speed; the LWR (Lighthill-Whitham-Richards) model as a PDE description of traffic dynamics; Signal control algorithms: fixed-time plans, actuated control, and the Webster formula for optimal cycle length as the classical control-theoretic approach to intersection management; Public transit networks: GTFS (General Transit Feed Specification) as the standard data schema for routes, stops, trips, and schedules; transit assignment and the stochastic user equilibrium problem; Shared mobility systems: bike-share and ride-hailing as queuing systems; station-based rebalancing as an optimization problem over demand prediction; Multimodal journey planning as a time-expanded graph shortest-path problem: integrating walking, transit, cycling, and ride-hailing into a unified routing engine.

\vspace{0.3cm}

\noindent\textbf{Unit 4} \\
\textit{Urban Infrastructure Management and Resilience:} Urban infrastructure systems as interdependent networks: power grids, water distribution, wastewater, telecommunications, and transportation as coupled graphs where failure in one propagates to others; Structural health monitoring (SHM): embedding vibration, strain, and tilt sensors in bridges and buildings; anomaly detection in sensor streams as the early-warning system for infrastructure degradation; Water distribution network modeling: pipe networks as hydraulic graphs governed by conservation of mass and energy; leak detection and pipe burst localization as inverse problems on sensor data; Smart grid fundamentals: the bidirectional power flow problem introduced by distributed renewable generation; demand response as a real-time optimization of load shedding under supply constraints; Resilience metrics for urban infrastructure: redundancy, robustness, rapidity, and resourcefulness as the 4R framework; cascading failure simulation using percolation theory on interdependent networks; Urban flood modeling: drainage network capacity, impervious surface coverage, and DEM-based inundation simulation as a coupled hydraulic-GIS computational problem.

\vspace{0.3cm}

\noindent\textbf{Unit 5} \\
\textit{Urban Governance, Open Data, and Civic Technology:} Open government data as a public good: data licensing (CC-BY, ODbL), machine-readable formats (CSV, GeoJSON, SPARQL endpoints), and the five-star open data maturity model; Urban dashboards and decision support systems: real-time KPI visualization for city operators as a data pipeline from sensor ingestion to front-end rendering; Participatory platforms and e-governance: 311 service request systems, participatory budgeting platforms, and fix-my-street applications as citizen-government data exchange interfaces; Algorithmic accountability in smart city systems: auditing predictive policing models, automated permit systems, and dynamic pricing algorithms for fairness and transparency; Smart city standards and frameworks: ISO 37120 (city indicators), ITU-T Y.4000 (IoT for smart cities), and the ISO/IEC 30145 smart city ICT reference framework as the interoperability layer enabling vendor-neutral city platforms; Privacy by design in urban sensing: anonymization of mobility traces, federated analytics for utility data, and the legal landscape of public surveillance under GDPR.

\newpage
\rhead{Curriculum Schema}
